\documentclass[11pt,a4paper]{article}
\usepackage[utf8]{inputenc}
\usepackage[T1]{fontenc}
\usepackage[spanish,es-minimal,es-noshorthands]{babel}
\AtBeginDocument{\renewcommand{\sin}{\operatorname{sen}}\renewcommand{\arcsin}{\operatorname{arcsen}}}
\usepackage{lmodern}
\usepackage[sfdefault,lining]{FiraSans}
\usepackage[a4paper,top=2.3cm,bottom=2.2cm,left=1.8cm,right=1.8cm,headheight=1.3cm,headsep=0.4cm,footskip=1cm]{geometry}
\usepackage{amsmath,amssymb}
\usepackage{icomma}
\usepackage{graphicx}
\usepackage[table]{xcolor}
\usepackage{tikz}
\usepackage{fontawesome5}
\usepackage{ccicons}
\usepackage{fancyhdr}
\usepackage{lastpage}
\usepackage{tasks}
\usepackage{needspace}
\usepackage[most]{tcolorbox}
\usepackage{array}
\usepackage[hidelinks]{hyperref}
\AtBeginDocument{\hypersetup{pdftitle={Soluciones de la ficha de la unidad 6: Introducción al álgebra. Ecuaciones de primer grado. 1º ESO},pdfsubject={Matemáticas 1º ESO},pdfkeywords={lenguaje algebraico, expresiones algebraicas, valor numérico, monomios, polinomios, propiedad distributiva, ecuaciones de primer grado, paréntesis, denominadores, problemas}}}

\definecolor{navy}{HTML}{1B2748}
\definecolor{accent}{HTML}{E8603A}
\definecolor{teal}{HTML}{1F9E90}
\definecolor{soft}{HTML}{EEF1F8}
\definecolor{muted}{HTML}{5D6685}

\setlength{\parindent}{0pt}
\setlength{\parskip}{3pt}
\settasks{label=\alph*), label-format=\bfseries\color{navy}, label-width=1.3em, item-indent=1.8em, column-sep=0.6em, after-item-skip=5pt, before-skip=3pt, after-skip=4pt}

% ---------- capçalera i peu ----------
\newcommand{\llicencia}{\href{https://creativecommons.org/licenses/by-nc-sa/4.0/deed.es}{\ccbyncsa}\ \ Licencia Creative Commons Atribución-NoComercial-CompartirIgual 4.0 Internacional (CC BY-NC-SA 4.0)}
\pagestyle{fancy}
\fancyhf{}
\fancyhead[L]{\raisebox{-0.35\height}{\includegraphics[height=0.95cm]{logo-gva.png}}}
\fancyhead[C]{\small\color{muted}Soluciones de la ficha · Unidad 6 · 1\textsuperscript{o} ESO}
\fancyhead[R]{\raisebox{-0.35\height}{\includegraphics[height=0.95cm]{logo-vermellar.png}}}
\fancyfoot[L]{\scriptsize\color{muted}\llicencia}
\fancyfoot[R]{\small\color{muted}\thepage}
\renewcommand{\headrulewidth}{0.4pt}
\renewcommand{\headrule}{\color{muted!40}\hrule height \headrulewidth}
\fancypagestyle{primera}{\fancyhf{}\renewcommand{\headrulewidth}{0pt}
  \fancyfoot[L]{\scriptsize\color{muted}\llicencia}
  \fancyfoot[R]{\small\color{muted}\thepage}}

% ---------- nivells i exercicis ----------
\newcommand{\estrelles}[1]{{\color{accent}\ifcase#1\or\faStar\faStar[regular]\faStar[regular]\or\faStar\faStar\faStar[regular]\or\faStar\faStar\faStar\fi}}
\newcounter{exer}
\newcommand{\ex}[1]{\par\needspace{4\baselineskip}\medskip\stepcounter{exer}\noindent{\color{navy}\bfseries\theexer.}\ \ {\small\estrelles{#1}}\ \ }

\newcommand{\seccio}[2]{\par\needspace{10\baselineskip}\bigskip
  \begin{tcolorbox}[colback=navy,colframe=navy,coltext=white,arc=2mm,boxrule=0pt,left=3mm,right=3mm,top=1.2mm,bottom=1.2mm]
  {\large\bfseries #1\quad #2}
  \end{tcolorbox}\vspace{-2pt}}

\newcolumntype{C}[1]{>{\centering\arraybackslash}p{#1}}

% quadre de recordatori (només a la fitxa d'exercicis)
\newcommand{\recorda}[1]{\par\smallskip\begin{tcolorbox}[colback=accent!7,colframe=accent!55,arc=2mm,boxrule=0.5pt,left=2.5mm,right=2.5mm,top=1mm,bottom=1mm]\small{\color{accent}\faLightbulb[regular]}\ \textbf{Recuerda:}\ #1\end{tcolorbox}}
\newcommand{\res}[1]{\boxed{#1}}\newcommand{\rest}[1]{\fbox{#1}}

% =====================================================================
% Capçalera FentMates de les fitxes: un color per a cada curs.
% Els colors són els de la pàgina inicial (LEVELS a assets/js/materials.js);
% el logo ve de _eines/logo/pdf/ (python3 _eines/logo/genera.py).
%
% Ús, al final del preàmbul de la unitat:
%   % =====================================================================
% Capçalera FentMates de les fitxes: un color per a cada curs.
% Els colors són els de la pàgina inicial (LEVELS a assets/js/materials.js);
% el logo ve de _eines/logo/pdf/ (python3 _eines/logo/genera.py).
%
% Ús, al final del preàmbul de la unitat:
%   % =====================================================================
% Capçalera FentMates de les fitxes: un color per a cada curs.
% Els colors són els de la pàgina inicial (LEVELS a assets/js/materials.js);
% el logo ve de _eines/logo/pdf/ (python3 _eines/logo/genera.py).
%
% Ús, al final del preàmbul de la unitat:
%   \input{../comu/fentmates.tex}
%   \fmcurs{1eso}
% i, al començament del document:
%   \capcalera{Matemàtiques}{1r ESO}{Fitxa d'exercicis · Unitat 2}{Divisibilitat}
% Opcional abans de \input: \def\fmneed{14} (espai mínim abans de cada secció).
% =====================================================================
\makeatletter
\def\fm@logodir{../../logo/pdf}
\providecommand{\fmneed}{10}
\def\fm@cursnom{}

% clau del curs, color i color del text damunt del color
\newcommand{\fm@curs}[3]{\@namedef{fm@color@#1}{#2}\@namedef{fm@text@#1}{#3}}
\fm@curs{1eso}{E6B84A}{navy}
\fm@curs{2eso}{3B82F6}{white}
\fm@curs{3eso}{FF7A50}{white}
\fm@curs{4eso-a}{A78BFA}{white}
\fm@curs{4eso-b}{6D28D9}{white}
\fm@curs{1batx-mi}{9B1C1C}{white}
\fm@curs{1batx-ccss}{EF6B6B}{white}
\fm@curs{2batx-mii}{1E6B47}{white}
\fm@curs{2batx-ccss}{2BB3A3}{white}

\newcommand{\fmcurs}[1]{%
  \@ifundefined{fm@color@#1}{\PackageError{fentmates}{Curs desconegut: #1}{}}{}%
  \edef\fm@tmp{\noexpand\definecolor{curs}{HTML}{\@nameuse{fm@color@#1}}%
               \noexpand\colorlet{curstext}{\@nameuse{fm@text@#1}}}\fm@tmp
  \colorlet{cursfosc}{curs!72!black}%
  \def\fm@logo{\fm@logodir/vertical-#1.pdf}%
  \def\fm@marca{\fm@logodir/marca-#1.pdf}}

% ---------- capçalera de la primera pàgina ----------
% #1 matèria (petit), #2 curs, #3 tipus de document i unitat, #4 títol
\newcommand{\capcalera}[4]{%
  \gdef\fm@cursnom{#2}%
  \par\noindent
  \begin{tcolorbox}[colback=curs,colframe=curs,boxrule=0pt,arc=2.2mm,width=3.5cm,height=2.15cm,
      valign=center,left=3mm,right=2mm,top=1mm,bottom=1mm,nobeforeafter,box align=center]
    \color{curstext}{\footnotesize #1}\par\vspace{1pt}{\fontsize{19}{21}\selectfont\bfseries #2\par}%
  \end{tcolorbox}\hspace{5mm}%
  \begin{minipage}[c]{\dimexpr\linewidth-3.5cm-5mm-2.9cm\relax}
    {\color{cursfosc}\small\bfseries\MakeUppercase{#3}}\par\vspace{2pt}
    {\color{navy}\fontsize{20}{23}\selectfont\bfseries #4\par}
  \end{minipage}\hfill
  \begin{minipage}[c]{2.6cm}\raggedleft\includegraphics[width=2.6cm]{\fm@logo}\end{minipage}\par
  \vspace{2pt}}

% ---------- barres de secció del color del curs ----------
\renewcommand{\seccio}[2]{\par\needspace{\fmneed\baselineskip}\bigskip
  \begin{tcolorbox}[colback=curs,colframe=curs,coltext=curstext,arc=2mm,boxrule=0pt,left=3mm,right=3mm,top=1.2mm,bottom=1.2mm]
  {\large\bfseries #1\quad #2}
  \end{tcolorbox}\vspace{-2pt}}

% ---------- estrelles de dificultat del color del curs (a Batxillerat no se'n posen) ----------
\@ifundefined{estrelles}{}{%
  \renewcommand{\estrelles}[1]{{\color{curs}\ifcase#1\or\faStar\faStar[regular]\faStar[regular]\or\faStar\faStar\faStar[regular]\or\faStar\faStar\faStar\fi}}}

% ---------- capçalera i peu de les pàgines ----------
% primera pàgina: logos institucionals al peu
\fancypagestyle{primera}{\fancyhf{}\renewcommand{\headrulewidth}{0pt}
  \fancyfoot[L]{\raisebox{-0.32\height}{\includegraphics[height=0.78cm]{logo-gva.png}}\hspace{5mm}%
    \raisebox{-0.32\height}{\includegraphics[height=0.78cm]{logo-vermellar.png}}}
  \fancyfoot[R]{\scriptsize\color{muted}\href{https://creativecommons.org/licenses/by-nc-sa/4.0/deed.ca}{\ccbyncsa}\ \ CC BY-NC-SA 4.0 · fentmates.org\qquad\small\thepage}}
% la resta: el símbol a l'esquerra i el curs a la dreta, amb el títol al centre com abans
\fancyhead[L]{\raisebox{-0.3\height}{\includegraphics[height=0.62cm]{\fm@marca}}}
\fancyhead[R]{\small\bfseries\color{cursfosc}\fm@cursnom}
\renewcommand{\headrule}{\color{curs!45}\hrule height \headrulewidth}
\makeatother

%   \fmcurs{1eso}
% i, al començament del document:
%   \capcalera{Matemàtiques}{1r ESO}{Fitxa d'exercicis · Unitat 2}{Divisibilitat}
% Opcional abans de \input: \def\fmneed{14} (espai mínim abans de cada secció).
% =====================================================================
\makeatletter
\def\fm@logodir{../../logo/pdf}
\providecommand{\fmneed}{10}
\def\fm@cursnom{}

% clau del curs, color i color del text damunt del color
\newcommand{\fm@curs}[3]{\@namedef{fm@color@#1}{#2}\@namedef{fm@text@#1}{#3}}
\fm@curs{1eso}{E6B84A}{navy}
\fm@curs{2eso}{3B82F6}{white}
\fm@curs{3eso}{FF7A50}{white}
\fm@curs{4eso-a}{A78BFA}{white}
\fm@curs{4eso-b}{6D28D9}{white}
\fm@curs{1batx-mi}{9B1C1C}{white}
\fm@curs{1batx-ccss}{EF6B6B}{white}
\fm@curs{2batx-mii}{1E6B47}{white}
\fm@curs{2batx-ccss}{2BB3A3}{white}

\newcommand{\fmcurs}[1]{%
  \@ifundefined{fm@color@#1}{\PackageError{fentmates}{Curs desconegut: #1}{}}{}%
  \edef\fm@tmp{\noexpand\definecolor{curs}{HTML}{\@nameuse{fm@color@#1}}%
               \noexpand\colorlet{curstext}{\@nameuse{fm@text@#1}}}\fm@tmp
  \colorlet{cursfosc}{curs!72!black}%
  \def\fm@logo{\fm@logodir/vertical-#1.pdf}%
  \def\fm@marca{\fm@logodir/marca-#1.pdf}}

% ---------- capçalera de la primera pàgina ----------
% #1 matèria (petit), #2 curs, #3 tipus de document i unitat, #4 títol
\newcommand{\capcalera}[4]{%
  \gdef\fm@cursnom{#2}%
  \par\noindent
  \begin{tcolorbox}[colback=curs,colframe=curs,boxrule=0pt,arc=2.2mm,width=3.5cm,height=2.15cm,
      valign=center,left=3mm,right=2mm,top=1mm,bottom=1mm,nobeforeafter,box align=center]
    \color{curstext}{\footnotesize #1}\par\vspace{1pt}{\fontsize{19}{21}\selectfont\bfseries #2\par}%
  \end{tcolorbox}\hspace{5mm}%
  \begin{minipage}[c]{\dimexpr\linewidth-3.5cm-5mm-2.9cm\relax}
    {\color{cursfosc}\small\bfseries\MakeUppercase{#3}}\par\vspace{2pt}
    {\color{navy}\fontsize{20}{23}\selectfont\bfseries #4\par}
  \end{minipage}\hfill
  \begin{minipage}[c]{2.6cm}\raggedleft\includegraphics[width=2.6cm]{\fm@logo}\end{minipage}\par
  \vspace{2pt}}

% ---------- barres de secció del color del curs ----------
\renewcommand{\seccio}[2]{\par\needspace{\fmneed\baselineskip}\bigskip
  \begin{tcolorbox}[colback=curs,colframe=curs,coltext=curstext,arc=2mm,boxrule=0pt,left=3mm,right=3mm,top=1.2mm,bottom=1.2mm]
  {\large\bfseries #1\quad #2}
  \end{tcolorbox}\vspace{-2pt}}

% ---------- estrelles de dificultat del color del curs (a Batxillerat no se'n posen) ----------
\@ifundefined{estrelles}{}{%
  \renewcommand{\estrelles}[1]{{\color{curs}\ifcase#1\or\faStar\faStar[regular]\faStar[regular]\or\faStar\faStar\faStar[regular]\or\faStar\faStar\faStar\fi}}}

% ---------- capçalera i peu de les pàgines ----------
% primera pàgina: logos institucionals al peu
\fancypagestyle{primera}{\fancyhf{}\renewcommand{\headrulewidth}{0pt}
  \fancyfoot[L]{\raisebox{-0.32\height}{\includegraphics[height=0.78cm]{logo-gva.png}}\hspace{5mm}%
    \raisebox{-0.32\height}{\includegraphics[height=0.78cm]{logo-vermellar.png}}}
  \fancyfoot[R]{\scriptsize\color{muted}\href{https://creativecommons.org/licenses/by-nc-sa/4.0/deed.ca}{\ccbyncsa}\ \ CC BY-NC-SA 4.0 · fentmates.org\qquad\small\thepage}}
% la resta: el símbol a l'esquerra i el curs a la dreta, amb el títol al centre com abans
\fancyhead[L]{\raisebox{-0.3\height}{\includegraphics[height=0.62cm]{\fm@marca}}}
\fancyhead[R]{\small\bfseries\color{cursfosc}\fm@cursnom}
\renewcommand{\headrule}{\color{curs!45}\hrule height \headrulewidth}
\makeatother

%   \fmcurs{1eso}
% i, al començament del document:
%   \capcalera{Matemàtiques}{1r ESO}{Fitxa d'exercicis · Unitat 2}{Divisibilitat}
% Opcional abans de \input: \def\fmneed{14} (espai mínim abans de cada secció).
% =====================================================================
\makeatletter
\def\fm@logodir{../../logo/pdf}
\providecommand{\fmneed}{10}
\def\fm@cursnom{}

% clau del curs, color i color del text damunt del color
\newcommand{\fm@curs}[3]{\@namedef{fm@color@#1}{#2}\@namedef{fm@text@#1}{#3}}
\fm@curs{1eso}{E6B84A}{navy}
\fm@curs{2eso}{3B82F6}{white}
\fm@curs{3eso}{FF7A50}{white}
\fm@curs{4eso-a}{A78BFA}{white}
\fm@curs{4eso-b}{6D28D9}{white}
\fm@curs{1batx-mi}{9B1C1C}{white}
\fm@curs{1batx-ccss}{EF6B6B}{white}
\fm@curs{2batx-mii}{1E6B47}{white}
\fm@curs{2batx-ccss}{2BB3A3}{white}

\newcommand{\fmcurs}[1]{%
  \@ifundefined{fm@color@#1}{\PackageError{fentmates}{Curs desconegut: #1}{}}{}%
  \edef\fm@tmp{\noexpand\definecolor{curs}{HTML}{\@nameuse{fm@color@#1}}%
               \noexpand\colorlet{curstext}{\@nameuse{fm@text@#1}}}\fm@tmp
  \colorlet{cursfosc}{curs!72!black}%
  \def\fm@logo{\fm@logodir/vertical-#1.pdf}%
  \def\fm@marca{\fm@logodir/marca-#1.pdf}}

% ---------- capçalera de la primera pàgina ----------
% #1 matèria (petit), #2 curs, #3 tipus de document i unitat, #4 títol
\newcommand{\capcalera}[4]{%
  \gdef\fm@cursnom{#2}%
  \par\noindent
  \begin{tcolorbox}[colback=curs,colframe=curs,boxrule=0pt,arc=2.2mm,width=3.5cm,height=2.15cm,
      valign=center,left=3mm,right=2mm,top=1mm,bottom=1mm,nobeforeafter,box align=center]
    \color{curstext}{\footnotesize #1}\par\vspace{1pt}{\fontsize{19}{21}\selectfont\bfseries #2\par}%
  \end{tcolorbox}\hspace{5mm}%
  \begin{minipage}[c]{\dimexpr\linewidth-3.5cm-5mm-2.9cm\relax}
    {\color{cursfosc}\small\bfseries\MakeUppercase{#3}}\par\vspace{2pt}
    {\color{navy}\fontsize{20}{23}\selectfont\bfseries #4\par}
  \end{minipage}\hfill
  \begin{minipage}[c]{2.6cm}\raggedleft\includegraphics[width=2.6cm]{\fm@logo}\end{minipage}\par
  \vspace{2pt}}

% ---------- barres de secció del color del curs ----------
\renewcommand{\seccio}[2]{\par\needspace{\fmneed\baselineskip}\bigskip
  \begin{tcolorbox}[colback=curs,colframe=curs,coltext=curstext,arc=2mm,boxrule=0pt,left=3mm,right=3mm,top=1.2mm,bottom=1.2mm]
  {\large\bfseries #1\quad #2}
  \end{tcolorbox}\vspace{-2pt}}

% ---------- estrelles de dificultat del color del curs (a Batxillerat no se'n posen) ----------
\@ifundefined{estrelles}{}{%
  \renewcommand{\estrelles}[1]{{\color{curs}\ifcase#1\or\faStar\faStar[regular]\faStar[regular]\or\faStar\faStar\faStar[regular]\or\faStar\faStar\faStar\fi}}}

% ---------- capçalera i peu de les pàgines ----------
% primera pàgina: logos institucionals al peu
\fancypagestyle{primera}{\fancyhf{}\renewcommand{\headrulewidth}{0pt}
  \fancyfoot[L]{\raisebox{-0.32\height}{\includegraphics[height=0.78cm]{logo-gva.png}}\hspace{5mm}%
    \raisebox{-0.32\height}{\includegraphics[height=0.78cm]{logo-vermellar.png}}}
  \fancyfoot[R]{\scriptsize\color{muted}\href{https://creativecommons.org/licenses/by-nc-sa/4.0/deed.ca}{\ccbyncsa}\ \ CC BY-NC-SA 4.0 · fentmates.org\qquad\small\thepage}}
% la resta: el símbol a l'esquerra i el curs a la dreta, amb el títol al centre com abans
\fancyhead[L]{\raisebox{-0.3\height}{\includegraphics[height=0.62cm]{\fm@marca}}}
\fancyhead[R]{\small\bfseries\color{cursfosc}\fm@cursnom}
\renewcommand{\headrule}{\color{curs!45}\hrule height \headrulewidth}
\makeatother

\fmcurs{1eso}

\begin{document}
\thispagestyle{primera}

\capcalera{Matemáticas}{1\textsuperscript{o} ESO}{Soluciones de la ficha · Unidad 6}{Introducción al álgebra.\\ Ecuaciones de primer grado}
\vspace{8pt}

\begin{tcolorbox}[colback=soft,colframe=soft,arc=2mm,left=3mm,right=3mm,top=1.5mm,bottom=1.5mm]
\small Intenta resolver cada ejercicio antes de mirar su solución. Si no te sale el mismo resultado, repasa los pasos o consulta los apuntes interactivos de la unidad: encontrarás herramientas que resuelven paso a paso cualquier ecuación de primer grado. Las soluciones fraccionarias se dan en forma irreducible.
\end{tcolorbox}

% =====================================================
\seccio{1}{Lenguaje algebraico}

\ex{1} \begin{tasks}(2)
  \task El doble de un número: $\res{2x}$
  \task Un número más 7: $\res{x+7}$
  \task El triple de un número menos 2: $\res{3x-2}$
  \task La mitad de un número: $\res{\dfrac{x}{2}}$
  \task El cuadrado de un número: $\res{x^2}$
  \task Un número menos su cuarta parte: $\res{x-\dfrac{x}{4}}$
  \task El consecutivo de un número: $\res{x+1}$
  \task El doble de un número más su triple: $\res{2x+3x}$
\end{tasks}

\ex{1} \begin{tasks}(2)
  \task $x+5$: \rest{Un número más 5}
  \task $3x$: \rest{El triple de un número}
  \task $x-1$: \rest{El anterior de un número ($x$ menos 1)}
  \task $\dfrac{x}{3}$: \rest{La tercera parte de un número}
  \task $2(x+4)$: \rest{El doble de la suma de un número y 4}
  \task $x^2-1$: \rest{El cuadrado de un número menos 1}
\end{tasks}\par (Hay otras maneras correctas de decirlo.)

\ex{2} \begin{tasks}(3)
  \task $\res{x+6}$
  \task $\res{x-4}$
  \task $\res{5x}$
  \task $\res{\dfrac{x}{2}}$
  \task $\res{x+5x}$
  \task $\res{5x+10}$
\end{tasks}

\ex{2} \begin{tasks}(2)
  \task $\res{x}$ y $\res{x+1}$
  \task $\res{2x}$ y $\res{2x+2}$
  \task $\res{2x+1}$
  \task $\res{x+(x+1)+(x+2)}$
  \task $\res{1,5x}$\,€
  \task $\res{20x}$ céntimos
\end{tasks}

\ex{2} \begin{tasks}(2)
  \task $\res{8n}$
  \task $\res{4c+2m}$
  \task $\res{60h}$
  \task $P=\res{4l}$, \ $A=\res{l^2}$
\end{tasks}

\ex{3} a) $5,\ 8,\ 11$; la figura 4 tiene $\res{14}$ (cada figura añade 3). \ b) $\res{3n+2}$ (2 puntos a la izquierda y 3 por cada columna). \ c) $3\cdot10+2=\res{32}$; \ $3\cdot50+2=\res{152}$.

\ex{3} \begin{tasks}(2)
  \task $\res{3x-\dfrac{x}{2}}$
  \task $\res{2\bigl(x+(x+1)\bigr)}$
  \task $\res{a^2+b^2}$
  \task $\res{(a+b)^2}$
\end{tasks}


% =====================================================
\seccio{2}{Expresiones algebraicas. Valor numérico}

\ex{1} \begin{tasks}(4)
  \task $4\cdot2-3=\res{5}$
  \task $4\cdot0-3=\res{-3}$
  \task $4\cdot(-1)-3=\res{-7}$
  \task $4\cdot5-3=\res{17}$
\end{tasks}

\ex{1} \begin{tasks}(2)
  \task $3^2+1=\res{10}$
  \task $2\cdot4+(-3)=\res{5}$
  \task $10-3\cdot(-2)=10+6=\res{16}$
  \task $\dfrac{8}{2}+5=\res{9}$
  \task $3\cdot(4-1)=\res{9}$
  \task $(-2)\cdot5-(-2)=-10+2=\res{-8}$
\end{tasks}

\ex{2} {\centering\renewcommand{\arraystretch}{1.35}\begin{tabular}{|c|C{0.9cm}|C{0.9cm}|C{0.9cm}|C{0.9cm}|C{0.9cm}|C{0.9cm}||}\hline $x$ & $-2$ & $-1$ & $0$ & $1$ & $2$ & $3$\\\hline $2x+1$ & $\res{-3}$ & $\res{-1}$ & $\res{1}$ & $\res{3}$ & $\res{5}$ & $\res{7}$\\\hline $x^2-x$ & $\res{6}$ & $\res{2}$ & $\res{0}$ & $\res{0}$ & $\res{2}$ & $\res{6}$\\\hline\end{tabular}\par}

\ex{2} \begin{tasks}(1)
  \task $A=\dfrac{10\cdot7}{2}=\res{35\text{\,cm}^2}$
  \task $P=2\cdot12+2\cdot4,5=24+9=\res{33\text{\,cm}}$
  \task $V=4^3=\res{64\text{\,cm}^3}$
\end{tasks}

\ex{2} $P=\res{5+0,1m}$. \quad $m=120$: \ $5+0,1\cdot120=5+12=\res{17\text{\,€}}$. \quad $m=45$: \ $5+0,1\cdot45=5+4,5=\res{9,50\text{\,€}}$.

\ex{3} Primera: $y$ aumenta de 4 en 4, así que empieza por $4x$; para $x=1$, $4-2=2$: \ $y=\res{4x-2}$. \quad Segunda: cada $y$ es el cuadrado de $x$ más 1: \ $y=\res{x^2+1}$.

\ex{3} a) $x=\res{6}$ ($3\cdot6+2=20$). \ b) $x=\res{7}$ o $x=\res{-7}$ ($(-7)^2=49$). \ c) $x=\res{4}$ (los dos valen 13).


% =====================================================
\seccio{3}{Monomios y polinomios}

\ex{1} {\centering\renewcommand{\arraystretch}{1.35}\begin{tabular}{|c|C{2.2cm}|C{2.2cm}|C{1.4cm}|}\hline Monomio & Coeficiente & Parte literal & Grado\\\hline $6x^3$ & $\res{6}$ & $\res{x^3}$ & $\res{3}$\\\hline $-a^2b$ & $\res{-1}$ & $\res{a^2b}$ & $\res{3}$\\\hline $xy$ & $\res{1}$ & $\res{xy}$ & $\res{2}$\\\hline $-7$ & $\res{-7}$ & $\res{\text{—}}$ & $\res{0}$\\\hline $\tfrac{3}{4}x^2y^2$ & $\res{\tfrac{3}{4}}$ & $\res{x^2y^2}$ & $\res{4}$\\\hline $-z$ & $\res{-1}$ & $\res{z}$ & $\res{1}$\\\hline\end{tabular}\par}

\ex{1} Son monomios: \ $\res{5x^2}$, \ $\res{-3ab}$, \ $\res{8}$, \ $\res{x^2y^3}$ \ y \ $\res{\dfrac{x}{5}}$ (es $\frac{1}{5}x$). \ No lo son: $x+1$ (es una suma), $\frac{4}{x}$ (la letra divide) y $2^x$ (la letra está en el exponente).

\ex{2} Con $x$: \ $\res{4x,\ x,\ -3x}$. \quad Con $x^2$: \ $\res{-2x^2,\ 5x^2}$. \quad Con $xy$: \ $\res{3xy,\ -xy}$. \quad Sin letra: \ $\res{7}$.

\ex{2} Por ejemplo: \ a) $\res{5x^3}$ y $\res{x^3}$; \ b) $\res{4x^5}$ (o $4x^2y^3$); \ c) $\res{9}$; \ d) $\res{-xy}$. \ (Hay muchas respuestas correctas.)

\ex{2} \begin{tasks}(2)
  \task $3x^2-5x+1$: \ \res{3} términos, grado \res{2}, término indep. $\res{1}$
  \task $7-x^3$: \ \res{2} términos, grado \res{3}, término indep. $\res{7}$
  \task $x^4+2x^2-x$: \ \res{3} términos, grado \res{4}, término indep. $\res{0}$
  \task $5x-8$: \ \res{2} términos, grado \res{1}, término indep. $\res{-8}$
\end{tasks}

\ex{3} a) $\res{-6x^2y}$ \quad b) $\res{4x^2y}$ \quad c) $\res{5x^2y}$ \ (porque $6x^2y-5x^2y=x^2y$).

\ex{3} a) $-3\cdot(-2)^2\cdot5=-3\cdot4\cdot5=\res{-60}$ \quad b) $(-1)^3-2\cdot(-1)+1=-1+2+1=\res{2}$ \quad c) $2\cdot\frac{1}{4}-\frac{3}{2}=\frac{1}{2}-\frac{3}{2}=\res{-1}$


% =====================================================
\seccio{4}{Operaciones con monomios. Reducir expresiones}

\ex{1} \begin{tasks}(4)
  \task $5x+3x=\res{8x}$
  \task $9a-4a=\res{5a}$
  \task $x+x+x=\res{3x}$
  \task $2x-7x=\res{-5x}$
  \task $4x^{2}+x^{2}=\res{5x^{2}}$
  \task $-3ab+8ab=\res{5ab}$
  \task $6y-y=\res{5y}$
  \task $-x-x=\res{-2x}$
\end{tasks}

\ex{1} \begin{tasks}(2)
  \task $3x+2+5x-1=\res{8x+1}$
  \task $7-2x+4-x=\res{-3x+11}$
  \task $x^{2}+3x-2x^{2}+x=\res{-x^{2}+4x}$
  \task $5a+2b-3a+b=\res{2a+3b}$
\end{tasks}

\ex{1} \begin{tasks}(3)
  \task $3\cdot 5x=\res{15x}$
  \task $2x\cdot 4x=\res{8x^{2}}$
  \task $(-3a)\cdot 2a=\res{-6a^{2}}$
  \task $x\cdot x^{2}=\res{x^{3}}$
  \task $5x\cdot (-2y)=\res{-10xy}$
  \task $(-4x^{2})\cdot (-x)=\res{4x^{3}}$
\end{tasks}

\ex{2} \begin{tasks}(3)
  \task $4(x+3)=\res{4x+12}$
  \task $2(5x-1)=\res{10x-2}$
  \task $-3(x+2)=\res{-3x-6}$
  \task $-(4x-7)=\res{-4x+7}$
  \task $x(x+5)=\res{x^{2}+5x}$
  \task $2x(3x-1)=\res{6x^{2}-2x}$
\end{tasks}

\ex{2} \begin{tasks}(2)
  \task $2(x+4)+3x=2x+8+3x=\res{5x+8}$
  \task $5x-(2x+3)=5x-2x-3=\res{3x-3}$
  \task $3(x-2)+2(x+5)=3x-6+2x+10=\res{5x+4}$
  \task $4(2x+1)-(x-6)=8x+4-x+6=\res{7x+10}$
\end{tasks}

\ex{2} \begin{tasks}(2)
  \task $3x+4$ \rest{no es reducible} (no son semejantes)
  \task $2(x+5)=\res{2x+10}$
  \task $-(x-3)=\res{-x+3}$
  \task $x\cdot x=\res{x^2}$ \ ($x+x=2x$)
\end{tasks}

\ex{3} \begin{tasks}(1)
  \task $3(2x-1)-2(x-4)+x=6x-3-2x+8+x=\res{5x+5}$
  \task $5-2(3x+1)+4(x-2)=5-6x-2+4x-8=\res{-2x-5}$
  \task $x(x+2)-3(x-1)=x^{2}+2x-3x+3=\res{x^{2}-x+3}$
\end{tasks}

\ex{3} a) $2(2x+1)+2(x-1)=4x+2+2x-2=\res{6x}$; para $x=4$: $\res{24}$. \quad b) $x+(x+3)+(2x-1)=\res{4x+2}$; para $x=4$: $\res{18}$.

\ex{3} Siempre obtienes \rest{el número que has pensado}: \ $\dfrac{2(x+5)-10}{2}=\dfrac{2x+10-10}{2}=\dfrac{2x}{2}=\res{x}$.


% =====================================================
\seccio{5}{Ecuaciones. Elementos y solución}

\ex{1} a) Miembros: $\res{5x-3}$ y $\res{2x+9}$; términos: $\res{5x,\ -3,\ 2x,\ 9}$; incógnita: $\res{x}$; grado \res{1}. \par b) Miembros: $\res{4(a+1)}$ y $\res{a-2}$; términos: $\res{4(a+1),\ a,\ -2}$; incógnita: $\res{a}$; grado \res{1}.

\ex{1} \begin{tasks}(2)
  \task $3\cdot3-1=8$: \rest{sí}
  \task $2\cdot(-4)+5=-3$ y $-4+1=-3$: \rest{sí}
  \task $5-2=3$ y $2\cdot2=4$: \rest{no}
  \task $\frac{6}{3}+1=3$ y $6-3=3$: \rest{sí}
\end{tasks}

\ex{1} \begin{tasks}(3)
  \task $x=\res{6}$
  \task $x=\res{11}$
  \task $x=\res{6}$
  \task $x=\res{14}$
  \task $x=\res{7}$
  \task $x=\res{5}$
\end{tasks}

\ex{2} \begin{tasks}(2)
  \task $x=\res{3}$: \ $9=9$
  \task $x=\res{-2}$: \ $-2=-2$
  \task $x=\res{0}$: \ $0=0$
  \task $x=\res{1}$: \ $2=2$
\end{tasks}

\ex{2} \begin{tasks}(2)
  \task \rest{Identidad} (al operar, $2x+2=2x+2$)
  \task \rest{Ecuación} sin solución ($1\neq2$)
  \task \rest{Identidad}
  \task \rest{Ecuación}, con solución $x=4$
\end{tasks}

\ex{3} Por ejemplo: \ a) $\res{2x=10}$; \ b) $\res{3(x+4)=6}$; \ c) $\res{x+1=x+3}$. \ (Hay muchas respuestas correctas.)


% =====================================================
\seccio{6}{Ecuaciones equivalentes. Resolución}

\ex{1} \begin{tasks}(2)
  \task $x=15-8\ \Rightarrow\ \res{x=7}$
  \task $x=-2+6\ \Rightarrow\ \res{x=4}$
  \task $x=4-9\ \Rightarrow\ \res{x=-5}$
  \task $x=-11+5\ \Rightarrow\ \res{x=-6}$
\end{tasks}

\ex{1} \begin{tasks}(3)
  \task $6x=42\ \Rightarrow\ x=\tfrac{42}{6}=\res{7}$
  \task $-4x=20\ \Rightarrow\ x=\tfrac{20}{-4}=\res{-5}$
  \task $9x=-27\ \Rightarrow\ x=\tfrac{-27}{9}=\res{-3}$
  \task $\text{m.c.m.}=5:\ \ \res{x=15}$
  \task $\text{m.c.m.}=4:\ \ \res{x=-12}$
  \task $-x=8\ \Rightarrow\ \res{x=-8}$
\end{tasks}

\ex{1} \begin{tasks}(2)
  \task $2x=11-3\ \Rightarrow\ 2x=8\ \Rightarrow\ x=\tfrac{8}{2}=\res{4}$
  \task $5x=21+4\ \Rightarrow\ 5x=25\ \Rightarrow\ x=\tfrac{25}{5}=\res{5}$
  \task $3x=1-10\ \Rightarrow\ 3x=-9\ \Rightarrow\ x=\tfrac{-9}{3}=\res{-3}$
  \task $-2x=13-7\ \Rightarrow\ -2x=6\ \Rightarrow\ x=\tfrac{6}{-2}=\res{-3}$
\end{tasks}

\ex{2} \begin{tasks}(2)
  \task $5x-3x=10-2\ \Rightarrow\ 2x=8\ \Rightarrow\ x=\tfrac{8}{2}=\res{4}$
  \task $8x-6x=1+5\ \Rightarrow\ 2x=6\ \Rightarrow\ x=\tfrac{6}{2}=\res{3}$
  \task $4x-x=-8-7\ \Rightarrow\ 3x=-15\ \Rightarrow\ x=\tfrac{-15}{3}=\res{-5}$
  \task $2x-5x=3+9\ \Rightarrow\ -3x=12\ \Rightarrow\ x=\tfrac{12}{-3}=\res{-4}$
\end{tasks}

\ex{2} \begin{tasks}(1)
  \task $3x-x-x=9-4\ \Rightarrow\ \res{x=5}$
  \task $-2x-3x=1-6-5\ \Rightarrow\ -5x=-10\ \Rightarrow\ x=\tfrac{-10}{-5}=\res{2}$
  \task $x+2x+3x-4x=10\ \Rightarrow\ 2x=10\ \Rightarrow\ x=\tfrac{10}{2}=\res{5}$
  \task $9x-4x-2x=12+3\ \Rightarrow\ 3x=15\ \Rightarrow\ x=\tfrac{15}{3}=\res{5}$
\end{tasks}

\ex{2} \begin{tasks}(2)
  \task $4x=10\ \Rightarrow\ x=\tfrac{10}{4}=\res{\tfrac{5}{2}}$
  \task $6x-2x=3-1\ \Rightarrow\ 4x=2\ \Rightarrow\ x=\tfrac{2}{4}=\res{\tfrac{1}{2}}$
  \task $3x=5+2\ \Rightarrow\ 3x=7\ \Rightarrow\ \res{x=\tfrac{7}{3}}$
  \task $10x=-4\ \Rightarrow\ x=\tfrac{-4}{10}=\res{-\tfrac{2}{5}}$
\end{tasks}

\ex{3} \begin{tasks}(1)
  \task $2x-7x+x=3-6-5\ \Rightarrow\ -4x=-8\ \Rightarrow\ x=\tfrac{-8}{-4}=\res{2}$
  \task $-3x+x-4x=-6-12\ \Rightarrow\ -6x=-18\ \Rightarrow\ x=\tfrac{-18}{-6}=\res{3}$
  \task $x-x=5-1\ \Rightarrow\ 0x=4\ \Rightarrow\ \res{\text{sin solución}}$
  \task $3x-x-2x=6-6\ \Rightarrow\ 0x=0\ \Rightarrow\ \res{\text{todas las } x \text{ son solución (identidad)}}$
\end{tasks}\par En c) llegamos a $0x=4$: ningún número multiplicado por 0 da 4. En d), los dos miembros son iguales: es una identidad.

\ex{3} $\res{3x+4=x+10}$; \ $3x-x=10-4\ \Rightarrow\ 2x=6\ \Rightarrow\ x=\tfrac{6}{2}=\res{3}$ \ Cada caja pesa \rest{3\,kg}.


% =====================================================
\seccio{7}{Ecuaciones con paréntesis}

\ex{1} \begin{tasks}(1)
  \task $3x+6=18\ \Rightarrow\ 3x=18-6\ \Rightarrow\ 3x=12\ \Rightarrow\ x=\tfrac{12}{3}=\res{4}$
  \task $5x-5=20\ \Rightarrow\ 5x=20+5\ \Rightarrow\ 5x=25\ \Rightarrow\ x=\tfrac{25}{5}=\res{5}$
  \task $2x+14=x+12\ \Rightarrow\ 2x-x=12-14\ \Rightarrow\ \res{x=-2}$
  \task $-2x+6=10\ \Rightarrow\ -2x=10-6\ \Rightarrow\ -2x=4\ \Rightarrow\ x=\tfrac{4}{-2}=\res{-2}$
\end{tasks}

\ex{2} \begin{tasks}(1)
  \task $4x-4=2x+6\ \Rightarrow\ 4x-2x=6+4\ \Rightarrow\ 2x=10\ \Rightarrow\ x=\tfrac{10}{2}=\res{5}$
  \task $6x+3=5x+7\ \Rightarrow\ 6x-5x=7-3\ \Rightarrow\ \res{x=4}$
  \task $-x+4=x\ \Rightarrow\ -x-x=-4\ \Rightarrow\ -2x=-4\ \Rightarrow\ x=\tfrac{-4}{-2}=\res{2}$
  \task $3x-5=10\ \Rightarrow\ 3x=10+5\ \Rightarrow\ 3x=15\ \Rightarrow\ x=\tfrac{15}{3}=\res{5}$
\end{tasks}

\ex{2} \begin{tasks}(1)
  \task $2x+10=16\ \Rightarrow\ 2x=16-10\ \Rightarrow\ 2x=6\ \Rightarrow\ x=\tfrac{6}{2}=\res{3}$
  \task $5x-8=12\ \Rightarrow\ 5x=12+8\ \Rightarrow\ 5x=20\ \Rightarrow\ x=\tfrac{20}{5}=\res{4}$
  \task $5x-10=5\ \Rightarrow\ 5x=5+10\ \Rightarrow\ 5x=15\ \Rightarrow\ x=\tfrac{15}{5}=\res{3}$
  \task $8x-12=5x+15\ \Rightarrow\ 8x-5x=15+12\ \Rightarrow\ 3x=27\ \Rightarrow\ x=\tfrac{27}{3}=\res{9}$
\end{tasks}

\ex{3} \begin{tasks}(1)
  \task $3x-8=x+4\ \Rightarrow\ 3x-x=4+8\ \Rightarrow\ 2x=12\ \Rightarrow\ x=\tfrac{12}{2}=\res{6}$
  \task $-2x+11=3x-19\ \Rightarrow\ -2x-3x=-19-11\ \Rightarrow\ -5x=-30\ \Rightarrow\ x=\tfrac{-30}{-5}=\res{6}$
  \task $x+10=2x+10\ \Rightarrow\ x-2x=10-10\ \Rightarrow\ -x=0\ \Rightarrow\ \res{x=0}$
  \task $x-8=x-8\ \Rightarrow\ x-x=-8+8\ \Rightarrow\ 0x=0\ \Rightarrow\ \res{\text{todas las } x \text{ son solución (identidad)}}$
\end{tasks}\par En d) los dos miembros quedan iguales: es una identidad y cualquier número es solución.

\ex{3} \begin{tasks}(1)
  \task $-x+10=-3x+18\ \Rightarrow\ -x+3x=18-10\ \Rightarrow\ 2x=8\ \Rightarrow\ x=\tfrac{8}{2}=\res{4}$
  \task $-2x+10=0\ \Rightarrow\ -2x=-10\ \Rightarrow\ x=\tfrac{-10}{-2}=\res{5}$
\end{tasks}


% =====================================================
\seccio{8}{Ecuaciones con denominadores}

\ex{1} \begin{tasks}(2)
  \task $\text{m.c.m.}=3:\ \ \res{x=15}$
  \task $\text{m.c.m.}=4:\ \ x+8=20\ \Rightarrow\ x=20-8\ \Rightarrow\ \res{x=12}$
  \task $\text{m.c.m.}=3:\ \ 2x=24\ \Rightarrow\ x=\tfrac{24}{2}=\res{12}$
  \task $\text{m.c.m.}=2:\ \ x-2=8\ \Rightarrow\ x=8+2\ \Rightarrow\ \res{x=10}$
\end{tasks}

\ex{2} \begin{tasks}(1)
  \task $\text{m.c.m.}=4:\ \ 2x+x=24\ \Rightarrow\ 3x=24\ \Rightarrow\ x=\tfrac{24}{3}=\res{8}$
  \task $\text{m.c.m.}=6:\ \ 2x-x=12\ \Rightarrow\ \res{x=12}$
  \task $\text{m.c.m.}=6:\ \ 3x+2x=60\ \Rightarrow\ 5x=60\ \Rightarrow\ x=\tfrac{60}{5}=\res{12}$
  \task $\text{m.c.m.}=5:\ \ 5x-x=40\ \Rightarrow\ 4x=40\ \Rightarrow\ x=\tfrac{40}{4}=\res{10}$
\end{tasks}

\ex{2} \begin{tasks}(1)
  \task $\text{m.c.m.}=2:\ \ x+3=10\ \Rightarrow\ x=10-3\ \Rightarrow\ \res{x=7}$
  \task $\text{m.c.m.}=3:\ \ x-1=3x-15\ \Rightarrow\ x-3x=-15+1\ \Rightarrow\ -2x=-14\ \Rightarrow\ x=\tfrac{-14}{-2}=\res{7}$
  \task $\text{m.c.m.}=5:\ \ 2x+1=15\ \Rightarrow\ 2x=15-1\ \Rightarrow\ 2x=14\ \Rightarrow\ x=\tfrac{14}{2}=\res{7}$
  \task $\text{m.c.m.}=10:\ \ 5x=2(x+6)\ \Rightarrow\ 5x=2x+12\ \Rightarrow\ 5x-2x=12\ \Rightarrow\ 3x=12\ \Rightarrow\ x=\tfrac{12}{3}=\res{4}$
\end{tasks}

\ex{2} \begin{tasks}(1)
  \task $\text{m.c.m.}=6:\ \ 3x+18=2x+30\ \Rightarrow\ 3x-2x=30-18\ \Rightarrow\ \res{x=12}$
  \task $\text{m.c.m.}=4:\ \ x+1=2(x-1)\ \Rightarrow\ x+1=2x-2\ \Rightarrow\ x-2x=-2-1\ \Rightarrow\ -x=-3\ \Rightarrow\ \res{x=3}$
\end{tasks}

\ex{3} \begin{tasks}(1)
  \task $\text{m.c.m.}=12:\ \ 4(x+2)-3(x-1)=24\ \Rightarrow\ x+11=24\ \Rightarrow\ x=24-11\ \Rightarrow\ \res{x=13}$
  \task $\text{m.c.m.}=3:\ \ 3x-(x+2)=18\ \Rightarrow\ 2x-2=18\ \Rightarrow\ 2x=18+2\ \Rightarrow\ 2x=20\ \Rightarrow\ x=\tfrac{20}{2}=\res{10}$
  \task $\text{m.c.m.}=4:\ \ 2(3x-1)-(x+5)=x\ \Rightarrow\ 5x-7=x\ \Rightarrow\ 5x-x=7\ \Rightarrow\ 4x=7\ \Rightarrow\ \res{x=\tfrac{7}{4}}$
\end{tasks}

\ex{3} \begin{tasks}(1)
  \task $\text{m.c.m.}=8:\ \ 4x+2x+x=112\ \Rightarrow\ 7x=112\ \Rightarrow\ x=\tfrac{112}{7}=\res{16}$
  \task $\text{m.c.m.}=6:\ \ 4(x-1)=3(x+4)\ \Rightarrow\ 4x-4=3x+12\ \Rightarrow\ 4x-3x=12+4\ \Rightarrow\ \res{x=16}$
\end{tasks}


% =====================================================
\seccio{9}{Resolución de problemas con ecuaciones}

\ex{1} $x+12=30\ \Rightarrow\ x=30-12\ \Rightarrow\ \res{x=18}$\par El número es el \rest{18}.

\ex{1} $3x-8=25\ \Rightarrow\ 3x=25+8\ \Rightarrow\ 3x=33\ \Rightarrow\ x=\tfrac{33}{3}=\res{11}$\par El número es el \rest{11}. ($33-8=25$)

\ex{1} $x+(x+1)=73\ \Rightarrow\ 2x+1=73\ \Rightarrow\ 2x=73-1\ \Rightarrow\ 2x=72\ \Rightarrow\ x=\tfrac{72}{2}=\res{36}$\par Los números son el \rest{36 y el 37}.

\ex{2} $x+(x+1)+(x+2)=159\ \Rightarrow\ 3x+3=159\ \Rightarrow\ 3x=159-3\ \Rightarrow\ 3x=156\ \Rightarrow\ x=\tfrac{156}{3}=\res{52}$\par Los números son \rest{52, 53 y 54}.

\ex{2} Hermano: $x$; Laura: $x+6$. \ $x+(x+6)=30\ \Rightarrow\ 2x+6=30\ \Rightarrow\ 2x=30-6\ \Rightarrow\ 2x=24\ \Rightarrow\ x=\tfrac{24}{2}=\res{12}$\par El hermano tiene \rest{12 años} y Laura, \rest{18 años}.

\ex{2} Trabajamos en céntimos ($4,50$\,€ $=450$ céntimos). Bolígrafo: $x$; libreta: $2x$. \ $x+2x=450\ \Rightarrow\ 3x=450\ \Rightarrow\ x=\tfrac{450}{3}=\res{150}$\par El bolígrafo cuesta \rest{150 céntimos $=1,50$\,€} y la libreta, \rest{3\,€}.

\ex{2} Altura: $x$; base: $x+5$. \ $2x+2(x+5)=46\ \Rightarrow\ 4x+10=46\ \Rightarrow\ 4x=46-10\ \Rightarrow\ 4x=36\ \Rightarrow\ x=\tfrac{36}{4}=\res{9}$\par La altura mide \rest{9\,cm} y la base, \rest{14\,cm}.

\ex{2} Primero: $x$; segundo: $2x$; tercero: $3x$. \ $x+2x+3x=90\ \Rightarrow\ 6x=90\ \Rightarrow\ x=\tfrac{90}{6}=\res{15}$\par Reciben \rest{15, 30 y 45 cromos}.

\ex{3} Dentro de $x$ años: madre $42+x$, hijo $10+x$. \ $42+x=3(10+x)\ \Rightarrow\ x+42=3x+30\ \Rightarrow\ x-3x=30-42\ \Rightarrow\ -2x=-12\ \Rightarrow\ x=\tfrac{-12}{-2}=\res{6}$\par Dentro de \rest{6 años} (48 y 16 años).

\ex{3} $\dfrac{x}{2}+\dfrac{x}{4}=21\ \Rightarrow\ \text{m.c.m.}=4:\ \ 2x+x=84\ \Rightarrow\ 3x=84\ \Rightarrow\ x=\tfrac{84}{3}=\res{28}$\par El número es el \rest{28}.

\ex{3} Chicas: $x$; chicos: $28-x$. \ $x+4=28-x\ \Rightarrow\ x+x=28-4\ \Rightarrow\ 2x=24\ \Rightarrow\ x=\tfrac{24}{2}=\res{12}$\par Hay \rest{12 chicas y 16 chicos}.

\ex{3} Coches: $x$; motos: $30-x$. \ $4x+2(30-x)=96\ \Rightarrow\ 2x+60=96\ \Rightarrow\ 2x=96-60\ \Rightarrow\ 2x=36\ \Rightarrow\ x=\tfrac{36}{2}=\res{18}$\par Hay \rest{18 coches y 12 motos}.

\ex{3} $x+2x+3x=180\ \Rightarrow\ 6x=180\ \Rightarrow\ x=\tfrac{180}{6}=\res{30}$\par Los ángulos miden \rest{$30^\circ$, $60^\circ$ y $90^\circ$}.


\end{document}
